% sduthesis-doc.tex — SDUThesis 使用手册
\documentclass[a4paper,12pt]{ctexbook}

% ===== 宏包 =====
\usepackage{geometry}
\geometry{left=2.5cm, right=2.5cm, top=2.5cm, bottom=2.5cm}
\usepackage{hyperref}
\hypersetup{colorlinks=true, linkcolor=blue!60!black, urlcolor=blue!60!black}
\usepackage{xcolor}
\usepackage{listings}
\usepackage{booktabs}
\usepackage{longtable}
\usepackage{enumitem}
\usepackage{tcolorbox}
\tcbuselibrary{listings, skins, breakable}
\usepackage{fancyhdr}
\usepackage{metalogo}

% ===== 页眉页脚 =====
\pagestyle{fancy}
\fancyhf{}
\fancyhead[L]{SDUThesis 使用手册}
\fancyhead[R]{\leftmark}
\fancyfoot[C]{\thepage}

% ===== 代码样式 =====
\lstdefinestyle{latex}{
  language=[LaTeX]TeX,
  basicstyle=\ttfamily\small,
  keywordstyle=\color{blue!70!black}\bfseries,
  commentstyle=\color{green!50!black}\itshape,
  stringstyle=\color{red!60!black},
  backgroundcolor=\color{gray!5},
  frame=single,
  rulecolor=\color{gray!30},
  breaklines=true,
  columns=fullflexible,
  morekeywords={SDUSetup, GetTitle, GetAuthor, GetStudentId, GetSchool,
    GetMajor, GetSupervisor, GetYear, GetMonth,
    makecoverpage, maketable, printbib,
    cnabstract, enabstract, cnkeywords, enkeywords,
    myappendix, acknowledgement,
    NewHook, UseHook, AddToHook},
}

% ===== 提示框 =====
\newtcolorbox{infobox}{
  colback=blue!5, colframe=blue!40!black,
  fonttitle=\bfseries, title=提示,
  breakable,
}
\newtcolorbox{warnbox}{
  colback=yellow!5, colframe=yellow!60!black,
  fonttitle=\bfseries, title=注意,
  breakable,
}

% ===== 元信息 =====
\title{\textbf{SDUThesis 使用手册} \\ \large 山东大学本科毕业论文 \LaTeX{} 模板}
\author{SDUThesis 维护团队}
\date{v1.4.0 \quad 2025年}

\begin{document}

\maketitle
\tableofcontents

% ============================================================
\chapter{简介}
% ============================================================

\section{什么是 SDUThesis}

SDUThesis 是山东大学本科毕业论文的 \LaTeX{} 模板，帮助毕业生专注于论文内容写作，而不用操心格式排版。

\section{特性}

\begin{itemize}[nosep]
  \item \textbf{一键配置} — \lstinline|\SDUSetup{}| 集中配置论文元信息
  \item \textbf{模块化架构} — 通过加载不同模块支持本科/盲审等多种场景
  \item \textbf{Hook 扩展} — 基于 \LaTeX3 Hook 机制，方便扩展
  \item \textbf{开箱即用} — Fandol 开源字体，无需额外安装
  \item \textbf{Overleaf 支持} — 一键在线编辑
  \item \textbf{CI/CD} — 自动编译 + 自动发布
\end{itemize}

\section{系统要求}

\begin{longtable}{ll}
\toprule
\textbf{项目} & \textbf{要求} \\
\midrule
\endfirsthead
\toprule
\textbf{项目} & \textbf{要求} \\
\midrule
\endhead
\bottomrule
\endfoot
TeX 发行版 & TeX Live 2024+ 或 MiKTeX \\
编译器 & XeLaTeX \\
参考文献 & Biber + Bib\LaTeX{} \\
中文字体 & Fandol（TeX Live 自带） \\
\end{longtable}

% ============================================================
\chapter{快速开始}
% ============================================================

\section{安装}

\subsection{方式一：克隆仓库}

\begin{lstlisting}[style=latex]
git clone https://github.com/h1s97x/sduthesis.git
\end{lstlisting}

\subsection{方式二：Overleaf}

点击 README 中的 \textbf{Open in Overleaf} 按钮，直接在线编辑。

\subsection{方式三：下载 Release}

从 GitHub Releases 页面下载最新版本的 zip 包。

\section{最小示例}

\begin{lstlisting}[style=latex]
\documentclass{sduthesis}

\SDUSetup{
  module     = undergraduate,
  title      = {基于XX的XX研究},
  author     = {张三},
  studentId  = {2020XXXXXX},
  school     = {计算机科学与技术学院},
  major      = {计算机科学与技术},
  supervisor = {李四},
  year       = {2025},
  month      = {6},
}

\begin{document}

\makecoverpage
\maketable

\begin{cnabstract}
这里是中文摘要。
\cnkeywords{关键词1，关键词2，关键词3}
\end{cnabstract}

\begin{enabstract}
This is the English abstract.
\enkeywords{keyword1, keyword2, keyword3}
\end{enabstract}

\mainmatter
\chapter{绪论}
论文正文从这里开始。

\backmatter
\printbib
\acknowledgement{感谢...}

\end{document}
\end{lstlisting}

\section{编译}

\begin{lstlisting}[language=bash]
# 使用 just（推荐）
just build

# 或手动编译
xelatex main
biber main
xelatex main
xelatex main
\end{lstlisting}

\begin{warnbox}
必须按 xelatex → biber → xelatex × 2 的顺序编译四次。
少编译一次都会导致引用显示为 ``??'' 或目录为空。
\end{warnbox}

% ============================================================
\chapter{配置参考}
% ============================================================

本章详细说明 \lstinline|\SDUSetup{}| 的所有配置项。

\section{基本信息}

\begin{longtable}{llcl}
\toprule
\textbf{键名} & \textbf{类型} & \textbf{必填} & \textbf{说明} \\
\midrule
\endfirsthead
\toprule
\textbf{键名} & \textbf{类型} & \textbf{必填} & \textbf{说明} \\
\midrule
\endhead
\bottomrule
\endfoot
\lstinline|title|      & 字符串 & 是 & 论文标题 \\
\lstinline|author|     & 字符串 & 是 & 作者姓名 \\
\lstinline|studentId|  & 字符串 & 是 & 学号 \\
\lstinline|school|     & 字符串 & 是 & 学院名称 \\
\lstinline|major|      & 字符串 & 是 & 专业名称 \\
\lstinline|supervisor| & 字符串 & 是 & 指导教师 \\
\lstinline|year|       & 字符串 & 是 & 毕业年份 \\
\lstinline|month|      & 字符串 & 是 & 毕业月份 \\
\end{longtable}

\section{样式配置}

\begin{longtable}{llcl}
\toprule
\textbf{键名} & \textbf{类型} & \textbf{必填} & \textbf{说明} \\
\midrule
\endfirsthead
\toprule
\textbf{键名} & \textbf{类型} & \textbf{必填} & \textbf{说明} \\
\midrule
\endhead
\bottomrule
\endfoot
\lstinline|lineSpread| & 数值 & 否 & 行距倍数，默认 1.25 \\
\lstinline|pageLeft|   & 数值 & 否 & 左边距（cm），默认 2.5 \\
\lstinline|pageRight|  & 数值 & 否 & 右边距（cm），默认 2.5 \\
\lstinline|pageTop|    & 数值 & 否 & 上边距（cm），默认 2.5 \\
\lstinline|pageBottom| & 数值 & 否 & 下边距（cm），默认 2.5 \\
\end{longtable}

\section{模块选择}

\begin{longtable}{lll}
\toprule
\textbf{值} & \textbf{说明} & \textbf{包含功能} \\
\midrule
\endfirsthead
\toprule
\textbf{值} & \textbf{说明} & \textbf{包含功能} \\
\midrule
\endhead
\bottomrule
\endfoot
\lstinline|undergraduate| & 本科论文（默认） & 封面、摘要、声明、页眉页脚 \\
\lstinline|blindreview|  & 盲审模式 & 基于本科，隐藏作者/导师信息 \\
\end{longtable}

\begin{lstlisting}[style=latex]
% 本科论文（默认）
\SDUSetup{module = undergraduate, ...}

% 盲审模式
\SDUSetup{module = blindreview, ...}
\end{lstlisting}

% ============================================================
\chapter{论文结构}
% ============================================================

\section{文档结构}

\begin{lstlisting}[style=latex]
\documentclass{sduthesis}
\SDUSetup{...}

\begin{document}
\makecoverpage          % 封面
\maketable              % 目录

% ----- 前言部分 -----
\begin{cnabstract}      % 中文摘要
  ...
  \cnkeywords{...}
\end{cnabstract}

\begin{enabstract}      % 英文摘要
  ...
  \enkeywords{...}
\end{enabstract}

% ----- 正文部分 -----
\mainmatter
\chapter{绪论}
\chapter{...}
\chapter{结论}

% ----- 后记部分 -----
\backmatter
\printbib               % 参考文献
\acknowledgement{...}   % 致谢
% \myappendix{...}      % 附录（可选）

\end{document}
\end{lstlisting}

\section{封面}

封面由 \lstinline|\makecoverpage| 自动生成，信息来源于 \lstinline|\SDUSetup{}|。

如需自定义校徽，可替换 \lstinline|assets/sdulogo.pdf| 文件。

\section{摘要}

\begin{lstlisting}[style=latex]
\begin{cnabstract}
  本文研究了...
  \cnkeywords{关键词1，关键词2，关键词3}
\end{cnabstract}

\begin{enabstract}
  This paper studies...
  \enkeywords{keyword1, keyword2, keyword3}
\end{enabstract}
\end{lstlisting}

\begin{infobox}
中文关键词之间用中文逗号（，）分隔，英文关键词之间用英文逗号（,）分隔。
\end{infobox}

\section{参考文献}

使用 Bib\LaTeX{} + Biber 管理参考文献，格式遵循 GB/T 7714-2015。

\begin{lstlisting}[style=latex]
% 在正文中引用
根据文献\citing{key1,key2}的研究...

% 在后记部分输出参考文献列表
\printbib
\end{lstlisting}

参考文献条目写在 \lstinline|data/ref/references.bib| 中：

\begin{lstlisting}[style=latex]
@article{key1,
  author  = {作者},
  title   = {标题},
  journal = {期刊名},
  year    = {2024},
  volume  = {1},
  number  = {1},
  pages   = {1--10},
}
\end{lstlisting}

\begin{warnbox}
\lstinline|\citing{}| 产生上标引用（如\textsuperscript{[1]}），
这是山大论文格式要求。如需行内引用，请使用
\lstinline|\cite{}| 或 \lstinline|\parencite{}|。
\end{warnbox}

\section{致谢}

\begin{lstlisting}[style=latex]
\acknowledgement{
  感谢我的导师...
}
\end{lstlisting}

\section{附录}

\begin{lstlisting}[style=latex]
\myappendix{
  \chapter{附录A标题}
  附录内容...
}
\end{lstlisting}

% ============================================================
\chapter{常见问题}
% ============================================================

\section{编译相关}

\begin{enumerate}
  \item \textbf{引用显示 ``??''} \\
    原因：编译次数不够。解决：执行完整编译流程 xelatex → biber → xelatex × 2。

  \item \textbf{目录为空} \\
    原因：同上，至少需要两次 xelatex 编译才能生成目录。

  \item \textbf{字体找不到} \\
    原因：系统缺少 Fandol 字体。解决：安装完整 TeX Live，
    或在 \lstinline|\SDUSetup{}| 中指定系统已有的字体。
\end{enumerate}

\section{Overleaf 相关}

\begin{enumerate}
  \item \textbf{Overleaf 上编译失败} \\
    确保编译器设置为 XeLaTeX（Menu → Compiler → XeLaTeX）。
\end{enumerate}

\section{自定义}

\begin{enumerate}
  \item \textbf{如何改行距？} \\
    \lstinline|\SDUSetup{lineSpread = 1.5}|

  \item \textbf{如何改页边距？} \\
    \lstinline|\SDUSetup{pageLeft = 3, pageRight = 3, pageTop = 3, pageBottom = 3}|

  \item \textbf{如何换校徽？} \\
    替换 \lstinline|assets/sdulogo.pdf| 文件即可。
\end{enumerate}

% ============================================================
\chapter{架构设计}
% ============================================================

本章面向开发者和贡献者，介绍 SDUThesis 的内部架构。

\section{分层架构}

\begin{lstlisting}[style=latex]
┌─────────────────────────────────────┐
│       用户层 (User Layer)           │
│   sdusetup.tex + data/ + main.tex  │
├─────────────────────────────────────┤
│       模块层 (Module Layer)         │
│   undergraduate.sty  blindreview.sty│
├─────────────────────────────────────┤
│       内核 (Core Engine)            │
│   sduthesis.cls                     │
│   ├── SDUSetup 引擎 (l3keys)       │
│   ├── Hook 系统                     │
│   ├── 字体加载器                    │
│   └── 页面引擎                      │
├─────────────────────────────────────┤
│       基础层 (Base Layer)           │
│   LaTeX3 + ctex + xeCJK + biblatex │
└─────────────────────────────────────┘
\end{lstlisting}

\section{Hook 系统}

内核定义了以下 Hook，模块通过 \lstinline|\AddToHook| 注入行为：

\begin{longtable}{ll}
\toprule
\textbf{Hook 名称} & \textbf{触发时机} \\
\midrule
\endfirsthead
\toprule
\textbf{Hook 名称} & \textbf{触发时机} \\
\midrule
\endhead
\bottomrule
\endfoot
\lstinline|sduthesis/after-setup|       & SDUSetup 之后 \\
\lstinline|sduthesis/before-cover|      & 封面前 \\
\lstinline|sduthesis/cover-style|       & 封面样式（模块覆盖） \\
\lstinline|sduthesis/frontmatter/begin| & 前言开始 \\
\lstinline|sduthesis/mainmatter/begin|  & 正文开始 \\
\lstinline|sduthesis/backmatter/begin|  & 后记开始 \\
\end{longtable}

\section{SDUSetup 实现原理}

SDUSetup 基于 \LaTeX3 l3keys 机制，分为三层：

\begin{enumerate}
  \item \textbf{存储层} — \lstinline|\tl_new:N \l__sdu_title_tl| 声明 token 变量
  \item \textbf{定义层} — \lstinline|\keys_define:nn { sdu } { title .tl_set:N = \l__sdu_title_tl }| 绑定键与变量
  \item \textbf{导出层} — \lstinline|\NewDocumentCommand \GetTitle {} { \l__sdu_title_tl }| 导出用户命令
\end{enumerate}

数据流：用户 \lstinline|\SDUSetup{title={...}}| → l3keys 写入变量 → \lstinline|\GetTitle| 读取变量。

\section{如何创建新模块}

详见 \texttt{doc/DEVELOP.md}。简要步骤：

\begin{enumerate}
  \item 创建 \lstinline|modules/sduthesis-<name>.sty|
  \item 声明包：\lstinline|\ProvidesPackage{sduthesis-<name>}|
  \item 添加模块特有配置键
  \item 通过 \lstinline|\AddToHook| 挂载行为
  \item 通过 \lstinline|\renewcommand| 覆盖排版
\end{enumerate}

% ============================================================
\appendix
\chapter{完整配置示例}
% ============================================================

\begin{lstlisting}[style=latex]
\documentclass{sduthesis}

\SDUSetup{
  % ----- 基本信息 -----
  module     = undergraduate,
  title      = {基于深度学习的图像识别算法研究},
  author     = {张三},
  studentId  = {2020012345},
  school     = {计算机科学与技术学院},
  major      = {计算机科学与技术},
  supervisor = {李四教授},
  year       = {2025},
  month      = {6},
  % ----- 样式配置（可选）-----
  lineSpread = 1.25,
  pageLeft   = 2.5,
  pageRight  = 2.5,
  pageTop    = 2.5,
  pageBottom = 2.5,
}

\begin{document}

\makecoverpage
\maketable

\begin{cnabstract}
  本文针对图像识别中的关键问题，提出了一种基于深度学习的新方法...
  \cnkeywords{深度学习，图像识别，卷积神经网络}
\end{cnabstract}

\begin{enabstract}
  This paper addresses key issues in image recognition
  and proposes a novel deep learning based method...
  \enkeywords{Deep Learning, Image Recognition, CNN}
\end{enabstract}

\mainmatter
\chapter{绪论}
\section{研究背景}
随着人工智能技术的快速发展...

\chapter{相关工作}
\section{深度学习概述}
深度学习是机器学习的一个重要分支\citing{lecun2015deep}...

\chapter{方法}
\chapter{实验}
\chapter{结论}

\backmatter
\printbib
\acknowledgement{感谢导师的悉心指导...}

\end{document}
\end{lstlisting}

\end{document}
